\chapter{Análisis del Panorama Actual del Lenguaje Ofensivo en Internet}

\section{Estadísticas Actuales}

Un estudio publicado este año analiza el desarrollo de la agresividad en la plataforma X (anteriormente conocida como Twitter) a partir de su adquisición por Elon Musk, quien implementó una política de moderación de contenido más flexible. Esta decisión podría estar relacionada con el aumento de los discursos de odio en la plataforma \cite{hickey2025x}.

El artículo fue publicado el 12 de febrero de 2025 y presenta un análisis que abarca desde principios de 2022 hasta el 9 de junio de 2023, utilizando como base publicaciones en inglés extraídas de la plataforma X. El estudio revela un incremento considerable de al menos un 50\% en los mensajes con contenido de odio tras la compra de la plataforma por parte de Elon Musk. Asimismo, se identificó una mayor interacción con este tipo de publicaciones: el promedio de "me gusta" en mensajes catalogados como discurso de odio creció en un 70\%, lo que sugiere no solo una mayor exposición sino también una preocupante aceptación por parte de los usuarios.

El análisis se apoya en distintos gráficos estadísticos que refuerzan sus conclusiones. Por ejemplo, la Figura 2.1 representa la variación semanal del número de publicaciones con lenguaje de odio antes y después del cambio de administración en la red. El gráfico muestra una tendencia al alza, con un punto crítico en abril de 2023. Este repunte se vincula con la aparición de una mujer transgénero en una campaña publicitaria, hecho que desató una ola de reacciones hostiles en línea y contribuyó a intensificar el volumen de mensajes ofensivos en la plataforma.









\vspace{0.5cm}

\begin{figure}[h!]
    \centering
    % \includegraphics[width=0.6\textwidth]{figura_2_1.png}
    \caption{Publicaciones de odio con el tiempo. Fuente: \cite{hickey2025x}}
\end{figure}

Por su parte, la Figura 2.2 ilustra cómo aumentó la participación de los usuarios en publicaciones clasificadas como discursos de odio. Específicamente, se evidencia un incremento notable en el promedio de “me gusta” que reciben estos tuits, con un aumento del 70\% tras los cambios en la política de moderación. Este dato sugiere no solo una mayor circulación de contenido dañino, sino también una creciente validación o aceptación por parte de los usuarios.












\vspace{0.5cm}

\begin{figure}[h!]
    \centering
    % \includegraphics[width=0.75\textwidth]{figura_2_2.png}
    \caption{Promedio de me gusta y republicaciones de (a) publicaciones de odio y (b) publicaciones de referencia antes y después de la adquisición de Musk. Fuente: \cite{hickey2025x}.}
\end{figure}

Finalmente, la Figura 2.3 revela cuáles categorías de odio aumentaron más visiblemente, revelando tendencias específicas en la expresión de agresividad. Esta figura es particularmente útil para identificar los grupos más afectados por la flexibilización en la moderación del contenido.

\begin{figure}[h!]
    \centering
    % \includegraphics[width=0.55\textwidth]{figura_2_3.png}
    \caption{Aumento de los discursos de odio para diferentes dimensiones del odio. Las líneas negras representan errores estándar. Fuente: \cite{hickey2025x}.}
\end{figure}

Alejándonos por un momento del entorno de las redes sociales, un estudio realizado por Preply los días 21 y 22 de enero de 2022 encuestó a 1,409 personas que juegan videojuegos, de las cuales el 58\% se identificó como hombre y el 42\% como mujer, con una edad promedio de 36 años \cite{preply2022games}. Esta encuesta arrojó resultados preocupantes: más del 90\% de los participantes afirmaron haber presenciado o sufrido algún tipo de abuso emocional o acoso mientras jugaban en línea. Además, los resultados muestran patrones preocupantes relacionados con el tipo de acoso más frecuente Tabla 2.1 y los videojuegos donde se experimenta con mayor intensidad este comportamiento Tabla 2.2.








\vspace{0.5cm}

\begin{table}[h!]
    \centering
    \begin{tabular}{l c}
        \hline
        \textbf{Tipo de acoso} & \textbf{Porcentaje} \\
        \hline
        Sobrenombres ofensivos & 53 \% \\
        Racismo & 44 \% \\
        Stalking & 37 \% \\
        Discurso de odio & 36 \% \\
        Insultos & 35 \% \\
        Lenguaje explícito & 34 \% \\
        Swatting & 26 \% \\
        Flaming & 26 \% \\
        Amenazas físicas & 22 \% \\
        Doxxing & 20 \% \\
        \hline
    \end{tabular}
    \caption{Tipos de acoso experimentados por jugadores en línea \textbf{Fuente:} \cite{preply2022games}}
\end{table}

\begin{table}[h!]
    \centering
    \begin{tabular}{l c}
        \hline
        \textbf{Tipo de videojuego} & \textbf{Porcentaje} \\
        \hline
        Shooters & 52 \% \\
        Multiplayer Online Battle Arena (MOBA) & 45 \% \\
        Estrategia en Tiempo Real & 27 \% \\
        Juego de rol & 23 \% \\
        Deportes & 21 \% \\
        \hline
    \end{tabular}
    \caption{Tipo de videojuegos con los mayores abusos \textbf{Fuente:} \cite{preply2022games}}
\end{table}

\section{Impacto del Lenguaje Ofensivo}

El portal PantallasAmigas presenta un análisis sobre los efectos de los comentarios negativos en redes sociales y otros espacios digitales, basado en un estudio realizado por la Fundación MAPFRE y la Universidad de Deusto \cite{mapfre2023redes}. El estudio encuestó a 2520 personas españolas (49 \% hombres y 51 \% mujeres), con el objetivo de conocer las percepciones sobre el impacto emocional que tienen las críticas, los insultos y otros tipos de interacciones negativas en internet. Entre los hallazgos más relevantes, destaca que las personas reaccionan con más ansiedad, miedo, tristeza y problemas de sueño tras recibir un comentario negativo. Además, las mujeres son las más afectadas en este tipo de situaciones, ya que llegan a sufrir hasta más del doble de inseguridad que los hombres (22\% en mujeres y 8\% en hombres). Estas conclusiones se ilustran mediante cuadros estadísticos incluidos en el informe \cite{pantallasamigas2024comentarios}.








\vspace{0.5cm}

\begin{figure}[h!]
    \centering
    % \includegraphics[width=0.78\textwidth]{figura_2_4.png}
    \caption{Uso de redes sociales según el sexo. Fuente: Adaptado de \cite{mapfre2023redes}.}
\end{figure}

\begin{figure}[h!]
    \centering
    % \includegraphics[width=0.78\textwidth]{figura_2_5.png}
    \caption{Temática de los comentarios negativos recibidos según el sexo. Fuente: Adaptado de \cite{mapfre2023redes}.}
\end{figure}














\vspace{0.5cm}

\begin{figure}[h!]
    \centering
    % \includegraphics[width=0.78\textwidth]{figura_2_6.png}
    \caption{Reacción emocional según el sexo. Fuente: Adaptado de \cite{mapfre2023redes}.}
\end{figure}

\section{Ejemplos de incidentes notorios con comentarios ofensivos}

La Figura 2.7 presenta una serie de publicaciones auténticas recogidas de distintas redes sociales, que evidencian la continua presencia de lenguaje ofensivo en estas plataformas, a pesar de contar con mecanismos de moderación automatizada y humana.

\begin{figure}[h!]
    \centering
    % \includegraphics[width=0.85\textwidth]{figura_2_7.png}
    \caption{Capturas de pantalla en redes sociales \textbf{[Elaboración Propia]}}
\end{figure}

Sin embargo, este problema no se limita únicamente a redes sociales. También se extiende a otros entornos digitales como los servicios de transmisión en vivo y los videojuegos en línea. En la Figu-







\vspace{0.5cm}

ra 2.8 se expone un caso registrado durante la emisión de una teleserie para todo público en la plataforma Kick, donde se visualizaron comentarios ofensivos en el chat en tiempo real. Por su parte, la Figura 2.9 muestra un ejemplo de contenido inapropiado identificado dentro del entorno competitivo del videojuego Valorant.

\begin{figure}[h!]
    \centering
    % \includegraphics[width=0.78\textwidth]{figura_2_8.png}
    \caption{Ofensividad en Kick \textbf{[Elaboración Propia]}}
\end{figure}

\begin{figure}[h!]
    \centering
    % \includegraphics[width=0.78\textwidth]{figura_2_9.png}
    \caption{Ofensividad en Valorant \textbf{[Elaboración Propia]}}
\end{figure}

Por otro lado, si bien persiste la circulación de mensajes ofensivos sin restricción, también se han documentado casos en los que los sistemas de moderación censuran erróneamente comentarios legítimos. En la plataforma TikTok, por ejemplo, se han reportado situaciones en las que publicaciones no ofensivas fueron retiradas o limitadas, como se detalla al final de esta sección.



\vspace{0.5cm}

\begin{figure}[h!]
    \centering
    % \includegraphics[width=0.82\textwidth]{figura_2_10.png}
    \caption{Ejemplos de comentarios censurados injustamente en la plataforma TikTok \textbf{[Elaboración Propia]}}
\end{figure}





\vspace{0.5cm}

\section{Mercado Global}

\subsection{Mercado Mundial de Moderación de Contenido Automatizado}

El mercado mundial de moderación de contenido automatizado ha experimentado un crecimiento significativo en los últimos años, impulsado por el aumento del contenido generado por usuarios en plataformas digitales como redes sociales, foros en línea y servicios de streaming \cite{bri2024automated}. El mercado se basa en tecnologías de Inteligencia Artificial (IA) y Machine Learning (ML) para identificar y eliminar contenido perjudicial de manera eficiente, incluso sin necesidad de moderadores humanos. El tamaño del mercado mundial de moderación de contenido automatizado se valoró aproximadamente en 0,83 billones de dólares en 2023 y se pronostica 3,4 en 2032, esto con una Tasa de Crecimiento Anual Compuesta (CAGR) de 18,5\% de 2023 a 2032.

\begin{figure}[h!]
    \centering
    % \includegraphics[width=0.75\textwidth]{figura_2_11.png}
    \caption{Tamaño estimado del mercado global de moderación de contenido automatizado. Adaptado y traducido de \cite{bri2024automated}.}
\end{figure}

\subsection{Mercado de Soluciones de Moderación de Contenido}

Según el mercado de soluciones de moderación de contenido, el mercado se divide en soluciones locales y basadas en la nube, donde las grandes empresas dominan el mercado debido a sus requisitos generalizados de control de contenido, pero las Pequeña y Mediana Empresa (PYME) están aumentando inesperadamente gracias a las ofertas de moderación \cite{bri2024solutions}.

El mercado de soluciones cuenta con una CAGR de aproximadamente 13,4\% de 2024 a 2032, este mercado se puede segmentar por aplicación o por tipo:



\vspace{0.5cm}

El mercado de soluciones por aplicación se puede clasificar en las redes sociales, el minorista de comercio electrónico y otros. Mientras que según el tipo, el mercado global se puede clasificar en servicios, software y plataforma, donde se pronostica una mayoría de servicios para el 2032, según la Figura 2.12.

\begin{figure}[h!]
    \centering
    % \includegraphics[width=0.65\textwidth]{figura_2_12.png}
    \caption{Mercado Global de soluciones de moderación de contenido pronosticado al año 2032. Adaptado y traducido de \cite{bri2024solutions}.}
\end{figure}

\subsection{Mercado Global de Servicios de Moderación de Contenido}

Según datos recientes de SNS Insider (2024), el mercado global de servicios de moderación de contenido fue valorado en USD 10.01 mil millones en 2023 y se proyecta que alcanzará los USD 30.75 mil millones para 2032, registrando una CAGR del 13.33\% en el período 2024--2032 \cite{sns2024forecast}. Este crecimiento refleja la creciente necesidad de herramientas de moderación automatizada, impulsada por el aumento del contenido generado por usuarios en plataformas digitales.

Se resume parte del informe en la Tabla 2.3 y la Tabla 2.4.










\vspace{0.5cm}

\begin{table}[h!]
    \centering
    \begin{tabular}{l c}
        \hline
        \textbf{Segmento} & \textbf{Participación} \\
        \hline
        Grandes Empresas & 66 \% \\
        Medios y Entretenimiento & 28 \% \\
        Imágenes & 47 \% \\
        América del Norte & 40 \% \\
        \hline
    \end{tabular}
    \caption{Participación en ingresos del mercado de servicios de moderación de contenido (2023) \textbf{[Elaboración Propia]}}
\end{table}

\begin{table}[h!]
    \centering
    \begin{tabular}{l c}
        \hline
        \textbf{Segmento} & \textbf{CAGR Estimado} \\
        \hline
        PYME & 15,14 \% \\
        Retail y Comercio Electrónico & 15,39 \% \\
        Video & 16,18 \% \\
        Asia-Pacífico & 15,04 \% \\
        \hline
    \end{tabular}
    \caption{CAGR estimado en servicios de moderación de contenido (2024--2032) \textbf{[Elaboración Propia]}}
\end{table}

Finalmente, se extrae una comparativa de las predicciones por los estudios de mercado.


\vspace{0.5cm}

\begin{figure}[h!]
    \centering
    % \includegraphics[width=0.78\textwidth]{figura_2_13.png}
    \caption{Proyección de los mercados globales al año 2032 \textbf{[Elaboración Propia]}}
\end{figure}





